% ------------------------------------------------------------
\subsection{Summary of Contributions}
\label{subsec:summary}
% ------------------------------------------------------------

This paper addresses a structural gap in the project portfolio
management literature: the absence of a formal, computationally
tractable framework for operational budget allocation across a
committed portfolio of FIDIC-governed projects under cash-flow
uncertainty.

Three contributions are delivered.

\textbf{C1 --- Gymnasium-compatible MDP environment.}
\texttt{PortfolioBudgetingEnv} is a fully parameterisable, open-source
Gymnasium environment encoding the complete FIDIC cash-flow cycle:
advance payment and progressive recovery, progress-weighted milestone
payments with configurable weights and timestep eligibility
constraints, retention withholding and release, treasury-draw interest
on project cash deficits, and a two-condition termination gate with a
project-specific cure-period counter. The sole stochastic primitive is
the budget-to-progress efficiency $\eta_{i,t}$, drawn each period from
a configurable distribution. All parameters --- schedule, contract,
and portfolio --- are specified through a configuration dictionary
whose entries are probability distributions, enabling domain
randomisation and calibration to any project-based sector without
code modification. Gymnasium API compliance is verified across 51
checks in 10 test sections. A SQLite audit ledger records every
allocation, cash-flow event, breach flag, and EVM index at each step,
making policy behaviour fully interpretable.

\textbf{C2 --- Formal MDP specification.}
The portfolio budget allocation problem is formalised as a
finite-horizon MDP $\mathcal{M} = (\mathcal{S}, \mathcal{A},
\mathcal{P}, R, \gamma, T)$ with state dimension $9N + 1$, a
$7N + 2$ dimensional observation vector derived from actionable
signals rather than raw EVM indices, a three-step action projection
guaranteeing feasibility for any policy output, and a reward function
equal to the discounted net portfolio cash flow. All structural
incentives --- completion, breach avoidance, milestone timing, and
treasury-draw deterrence --- emerge from the financial mechanics
without auxiliary penalty terms. The Markov property holds exactly:
the full portfolio state $S_t$ is a sufficient statistic for the
observable history.

\textbf{C3 --- Benchmark evaluation and quantified optimality gap.}
Three deterministic greedy heuristics (Reach-Plan, Catchup-T,
Catchup-Next-T) are evaluated against a formally verified
full-information optimum---confirmed optimal by exhaustive enumeration
over all feasible allocation sequences under fixed efficiency
realisations---across four benchmark configurations of
increasing structural complexity, with fixed efficiency sequences
ensuring exact reproducibility. The best-heuristic optimality gap
grows from $3.40$ NPV units on the isolated single-project atom to
$72.15$ units on the three-project coordinated portfolio --- a
$21\times$ increase. On the most complex configuration, every greedy
heuristic produces a loss while the full-information optimum remains
profitable, establishing both the validity of the environment
mechanics and a concrete performance target for future learned agents.

% ------------------------------------------------------------
\subsection{Implications}
\label{subsec:implications}
% ------------------------------------------------------------

\paragraph{For sequential decision-making research.}
The environment provides a problem instance that is simultaneously
realistic (FIDIC mechanics calibrated to empirical construction
practice), formally specified (closed-form MDP with exact Markov
property), and computationally tractable (polynomial step cost,
episode length $T \sim 20$--$30$ periods for representative
portfolios). The intractability of exact stochastic optimisation ---
$K^{NT}$ scenario-tree leaves, $K^{72} \approx 10^{50}$ for a
moderate three-project instance --- makes this a well-motivated
benchmark for RL and approximate dynamic programming methods. The
full-information optimum provides an exact upper bound
(Proposition~\ref{prop:upper_bound}) against which any policy can
be evaluated without requiring a solver.

\paragraph{For project management practice.}
The experiment results show that no fixed rule performs consistently
well: the optimal policy is adaptive to the efficiency trajectory,
the milestone payment structure, the advance trigger timing, and the
phasing of concurrent projects simultaneously. For a portfolio manager,
this implies that rule-of-thumb allocation --- even well-calibrated
rules such as Catchup-T, which tracks the contractual tolerance floor
--- leaves substantial value on the table as portfolio complexity
grows. The audit ledger provides a deployment-ready interpretability
layer: every allocation, payment, and breach event is recorded per
period, making the environment suitable as a decision-support tool
even before a learned policy is trained.

\paragraph{For domain randomisation and generalisation.}
Each parameter in the configuration dictionary can be specified as a
probability distribution, and $N$ (portfolio size) can itself be a
configurable random variable. This design supports episode-level
domain randomisation \citep{Tobin2017} natively, without any
modification to the environment code. A policy trained across a
distribution of configurations is exposed to diverse advance
structures, efficiency profiles, milestone weights, and portfolio
sizes within each training run, a prerequisite for generalisation
across the range of contract types and market conditions encountered
in practice.

% ------------------------------------------------------------
\subsection{Limitations}
\label{subsec:limitations}
% ------------------------------------------------------------

The present formulation makes three deliberate simplifications, each
of which is a natural boundary for future extension.

\textbf{Independent projects.} Projects interact only through the
shared budget $B_t$. Shared workforce or equipment pools, technology
spillovers, and client-side interdependencies are not modelled.
Extending to resource-coupled projects would require augmenting the
state space with resource inventories and adding coupling constraints
to the action projection.

\textbf{Milestone payments only.} The payment structure follows FIDIC
milestone-certification mechanics. Time-based, cost-plus, and hybrid
contract forms are not currently supported, though the configuration
dictionary architecture accommodates them through additional milestone
profile parameters.

\textbf{Single stochastic primitive.} The budget-to-progress
efficiency $\eta_{i,t}$ absorbs all uncertainty. Scope changes,
external price shocks, and client payment delays are exogenous. This
restriction cleanly identifies the allocation problem structure; a
multi-source uncertainty extension would require additional state
components to maintain the Markov property.

% ------------------------------------------------------------
\subsection{Future Work}
\label{subsec:future}
% ------------------------------------------------------------

The environment and benchmark established here define a clear research
agenda.

\textbf{RL agent training.}
The immediate next step is to train a policy --- PPO \citep{Schulman2017}
is a natural starting point given its stability in continuous action
spaces --- on a domain-randomised distribution of configurations and
evaluate it against the benchmark configurations of
Section~\ref{sec:evaluation}. The full-information gap of
Figure~\ref{fig:gap_plot} provides the evaluation target; a trained
policy that closes a substantial fraction of the TC-3XMAIF-R gap
($72.15$ units) would constitute a meaningful improvement over all
greedy baselines.

\textbf{LLM-augmented policies.}
The observation vector $o_t$ is structured, low-dimensional
($7N + 2$ components), and semantically interpretable --- each
component has a direct financial or contractual meaning. This makes
it a natural input for large language model-based reasoning agents
integrated via LangChain or similar frameworks, with the benchmark
heuristics serving as reward signals for in-context action scoring.
The primary challenge for such an approach is generalisation across
configurations: a policy that overfits to the four benchmark
efficiency sequences does not transfer to novel portfolios. Generalisation
configurations designed to inject controlled uncertainty --- varied
S-curve shapes, heterogeneous milestone weights, mixed advance
structures --- are the appropriate next step for evaluating policy
robustness.

\textbf{Formal algorithmic upper bound.}
The exhaustively enumerated $Z^*$ is verified optimal for the four
fixed-efficiency configurations but does not extend to stochastic settings. A Linear
Programming relaxation of the deterministic oracle problem
\eqref{eq:ub_obj}, or a Dynamic Programming solution for small
instances, would provide a tighter and algorithmically verifiable
bound against which trained policies can be compared.

\textbf{Real-world deployment.}
Integration with Oracle Primavera P6 or equivalent portfolio management
information systems would enable deployment as a decision-support
tool, with live EVM exports providing the state vector and the
environment serving as a scenario simulator for prospective allocation
decisions. The SQLite audit ledger is structurally compatible with
standard EVM reporting schemas and requires no additional instrumentation.